\chapter{核心功能实现}

\section{基于角色的访问控制与身份认证实现}

多角色协同系统的首要问题是如何在统一平台中实现“同库共用而不同权访问”。本文系统采用 JWT 与 Spring Security 结合的方式完成统一身份认证，并在接口级和服务级同时实现角色约束。其核心处理流程可概括为：用户提交凭证 $\rightarrow$ 后端完成认证 $\rightarrow$ 生成带角色信息的 JWT $\rightarrow$ 前端携带令牌访问业务接口 $\rightarrow$ 后端过滤器解析令牌并恢复安全上下文 $\rightarrow$ 控制器和服务层继续执行角色与业务状态校验。

在工程实现上，安全控制链具有以下几个层次。第一层为登录认证层，主要负责账号、密码和启用状态校验，并生成具有有效期的令牌。第二层为路径授权层，即在安全配置中为不同接口路径指定允许访问的角色集合。第三层为业务授权层，即使用户具有接口访问权限，若目标业务对象当前状态不允许该角色执行某动作，系统仍会在服务层拒绝请求。例如，生产管理员不能在订单尚未进入生产阶段时提交生产完工，仓库管理员不能在订单未通过销售审核时执行发货，质检管理员也不能对非待检批次重复提交质检结果。

为进一步说明权限控制机制，可将系统的访问控制规则抽象为如下集合约束：设用户集合为 $U$，角色集合为 $R$，接口集合为 $A$，则存在映射关系

\begin{equation}
\phi: U \rightarrow 2^{R}
\end{equation}

表示每个用户可持有的角色集合；同时存在授权映射

\begin{equation}
\psi: A \rightarrow 2^{R}
\end{equation}

表示每个接口允许访问的角色集合。当用户 $u$ 访问接口 $a$ 时，只有当

\begin{equation}
\phi(u) \cap \psi(a) \neq \varnothing
\end{equation}

时，请求才有机会继续进入业务处理阶段；之后还需满足目标对象状态条件，才能最终执行成功。该模型说明，角色授权只是访问控制的第一步，而业务状态校验则保证了流程正确性。

在本系统中，该机制产生了明显的工程价值。其一，系统管理员、顾客、销售管理员、仓库管理员、生产管理员、采购管理员、供应商和质检管理员可以在同一平台上完成不同任务，而无须为每类用户维护独立系统。其二，权限控制与工作流状态结合后，能够有效避免“角色正确但时机错误”的业务异常。其三，JWT 机制降低了服务端状态维护成本，更适合本文采用的前后端分离结构。因此，认证授权模块不仅是安全组件，也是整个系统流程可控性的基础前提。

\section{订单—仓储—生产—质检协同流程实现}

订单到质检的协同链路是本文系统最核心的业务闭环，其本质是围绕订单对象组织库存、生产和质检三类资源的连续协同。项目实现中，该闭环主要由 \texttt{OrderWorkflowService} 和 \texttt{QualityService} 共同驱动，其状态推进并非由前端自由跳转，而是严格受限于订单当前状态、库存结果和质检结论。

从状态设计看，订单主要经历“待销售审核”“待仓库核查”“已接单”“生产中”“待质检”“待生产入库”“已发货”等阶段。其中，销售管理员将订单推进到仓库核查；仓库管理员根据库存检查结果决定进入已接单或生产中；生产管理员完工后使订单进入待质检；质检管理员依据结果将订单退回生产或推进为待生产入库；仓库管理员确认合格批次入库后，订单重新具备发货条件。这一设计保证了每次状态迁移都对应明确的责任主体。

为便于说明订单协同机制，可将其状态转移关系抽象为表~\ref{tab:order-status-flow}。

\begin{table}[htbp]
\centering
\footnotesize
\caption{订单—仓储—生产—质检协同的主要状态迁移}
\label{tab:order-status-flow}
\begin{tabular}{|c|c|p{5.6cm}|}
\hline
当前状态 & 下一状态 & 触发条件 \\\hline
待销售审核 & 待仓库核查 & 销售管理员审核通过 \\\hline
待销售审核 & 已拒绝 & 销售管理员拒绝并填写理由 \\\hline
待仓库核查 & 已接单 & 仓库核查库存充足并完成预留 \\\hline
待仓库核查 & 生产中 & 仓库核查发现库存不足，生成生产计划 \\\hline
生产中 & 待质检 & 生产管理员提交完工，系统生成待检批次 \\\hline
待质检 & 生产中 & 质检不合格，订单回退至生产阶段 \\\hline
待质检 & 待生产入库 & 所有关联批次质检合格 \\\hline
待生产入库 & 已接单 & 仓库确认生产入库成功并完成库存预留 \\\hline
已接单 & 已发货 & 仓库执行发货并扣减库存 \\\hline
\end{tabular}
\end{table}

在库存充足场景下，订单工作流相对简洁。仓库核查环节通过遍历订单明细并计算产品可用库存，若所有产品均满足需求，则系统调用预留库存逻辑，为每个库存条目增加 \texttt{reservedQuantity}，从而为后续发货预留资源。发货时，系统优先消耗预留数量，并为每个实际出库动作生成库存流水，这使得“预留”与“实出”在数据层面相互独立又逻辑连贯。

在库存不足场景下，系统会针对缺口产品自动创建生产计划。若设订单中第 $i$ 个产品的需求量为 $d_i$，当前系统可用库存为 $a_i$，则其缺口量 $s_i$ 可表示为

\begin{equation}
 s_i = \max(0, d_i - a_i)
\end{equation}

当 $s_i > 0$ 时，系统即为该产品生成一条生产计划，并将订单状态推进为“生产中”。这一处理机制说明，生产模块不是独立启动的，而是直接由订单履约缺口触发。

生产完工后，系统不会立即增加库存，而是先通过 \texttt{prepareCompletedPlanBatchesForQuality} 逻辑生成待检批次。批次以订单号、产品和完工数量为核心，记录生产责任人与质检状态。随后，质检模块对批次执行检验：若结果为不合格，系统向对应生产管理员发送返工通知，并通过 \texttt{updateRelatedOrderStatus} 将订单退回“生产中”；若结果为合格，且订单对应的所有批次均已合格，则订单进入“待生产入库”，仓库方可在后续完成生产入库确认。由此可见，质检并非附加步骤，而是决定订单能否重新进入仓储环节的关键门禁。

该闭环的工程价值主要体现在三个方面。首先，订单流程不再依赖人工催促，而是通过系统状态自动推进；其次，库存、生产和质检之间的因果关系通过实体和状态得以显式表达；最后，系统能够在每个关键节点保留操作记录，为后续追责和质量分析提供数据支持。

\section{采购—供应商—仓库协同流程实现}

与订单履约闭环相对应，原材料采购形成了另一条跨主体协同链路。该链路的显著特点在于其不仅涉及企业内部采购管理员和仓库管理员，还涉及外部供应商账号，因此对身份隔离、状态同步和业务匹配校验的要求更高。项目中，该流程主要由 \texttt{ProcurementWorkflowService} 实现。

采购流程开始于采购管理员创建采购单。系统首先校验采购明细是否为原材料类型产品，再校验原材料档案中是否已经维护优选供应商信息，最后检验所选供应商是否与档案中的供应商一致。若上述条件全部满足，系统自动生成采购单号并将采购单状态设为“待供应商确认”。这一设计不仅减少了人工输入采购单号带来的错误，也将原材料档案约束直接前置到采购创建环节。

采购单的状态在后续流程中依次可能经历“待供应商确认”“供应商已接单”“供应商已拒绝”“供应商已发货”“待仓库收货”“已入库”等阶段。其主要状态转换关系如表~\ref{tab:purchase-status-flow} 所示。

\begin{table}[htbp]
\centering
\footnotesize
\caption{采购—供应商—仓库协同的主要状态迁移}
\label{tab:purchase-status-flow}
\begin{tabular}{|c|c|p{5.6cm}|}
\hline
当前状态 & 下一状态 & 触发条件 \\\hline
待供应商确认 & 供应商已接单 & 供应商选择接受采购单 \\\hline
待供应商确认 & 供应商已拒绝 & 供应商拒绝采购单并给出说明 \\\hline
供应商已接单 & 供应商已发货 & 供应商提交发货信息 \\\hline
供应商已发货 & 待仓库收货 & 采购管理员通知仓库收货 \\\hline
待仓库收货 & 已入库 & 仓库完成收货并自动写入原材料库存 \\\hline
\end{tabular}
\end{table}

在该流程中，供应商与原材料档案之间的一致性校验尤为关键。若设原材料 $m$ 的优选供应商标识为 $P_m$，当前采购单所选供应商标识为 $S$，则系统仅在二者满足相关匹配关系时允许创建采购单，即可表述为

\begin{equation}
\operatorname{match}(P_m, S) = \text{true}
\end{equation}

否则采购单创建请求将被拒绝。项目中，匹配逻辑综合考虑了供应商姓名、编码和邮箱等信息，从而适配不同档案维护方式。这一机制有效保证了采购流程不会偏离原材料主数据约束。

在仓库收货阶段，系统会遍历采购明细，为每个原材料查找或创建对应的库存记录，并将数量累加到 \texttt{InventoryItem} 中，同时写入入库流水 \texttt{StockTransaction}。因此，采购流程不是在供应商发货时结束，而是以仓库确认并完成数据库层库存更新为真正闭环。该设计将“下单”“发货”“收货”和“入库”四个概念区分开来，提升了业务语义精度。

在工程意义上，该协同流程实现了企业内部采购行为与外部供应商响应之间的结构化对接，降低了传统依赖电话、聊天工具或纸质单据导致的信息遗漏风险。同时，系统通过采购单状态和实时通知将不同主体串联起来，使企业能够更清晰地掌握原材料供应链执行情况。

\section{库存预警与周计划自动生成算法实现}

库存预警与周计划模块是本文系统从事务处理延伸至业务决策支持的重要体现。其实现目标并非构建复杂的智能优化模型，而是在现有业务数据基础上给出具备可解释性的缺口识别和计划建议。项目中的相关逻辑主要分布在 \texttt{InventoryAlertController} 与 \texttt{WeeklyPlanningService} 中。

\subsection{库存预警判定机制}

库存预警首先面向成品和原材料分别计算短缺量。对于某一产品 $p$，设当前总库存为 $Q_p$，已预留库存为 $R_p$，则可用库存 $A_p$ 为

\begin{equation}
A_p = \max(0, Q_p - R_p)
\end{equation}

进一步设安全库存为 $S_p$，若该产品属于成品，则系统还需要考虑在制量 $P_p$；若属于原材料，则需考虑在途采购量 $T_p$。项目中对短缺量和推荐动作量的计算可分别表述为

\begin{equation}
Shortage_p = \max(0, S_p - A_p)
\end{equation}

\begin{equation}
Recommend_p = \max(0, S_p - A_p - Pipeline_p)
\end{equation}

其中，当 $p$ 为成品时，$Pipeline_p = P_p$；当 $p$ 为原材料时，$Pipeline_p = T_p$。若 \texttt{Recommend\_p} 大于零，系统便生成预警项。对成品而言，推荐动作是创建生产计划；对原材料而言，推荐动作是通知采购管理员创建采购申请。

该判定机制的优势在于其直接基于项目当前已有数据即可计算，无需额外引入复杂训练数据或预测模型。同时，通过区分“短缺量”和“推荐动作量”，系统既能反映当前安全库存缺口，也能考虑在制和在途因素，避免重复下达计划或采购。

\subsection{生产周计划生成机制}

生产周计划的生成逻辑与当前项目实现保持一致。设某成品 $p$ 上周完工量为 $H_p$，当前未完成订单需求为 $D_p$，当前可用成品库存为 $A_p$，增长系数为 $\alpha = 1.10$，则系统首先计算基线产量

\begin{equation}
B_p = \max(\alpha H_p, D_p)
\end{equation}

然后计算建议生产量

\begin{equation}
Q_p^{prod} = \max(0, B_p - A_p)
\end{equation}

其中，$Q_p^{prod}$ 即写入生产周计划明细的建议生产量。项目中还会在明细中保存上周完工量、未完成需求量、当前库存量、增长系数和建议原因，便于业务人员理解计划来源。测试样例显示，在参考日期为 2026 年 4 月 15 日的场景下，系统可基于上周完工量 10、当前需求 8、库存 3 计算出基线 11 和建议生产量 8，这说明该规则能够与业务数据形成稳定对应关系。

\subsection{采购周计划生成机制}

采购周计划进一步以生产周计划和 BOM 关系为输入。设原材料 $m$ 由 BOM 换算得到的需求量为 $B_m$，上周采购量为 $H_m$，当前可用原材料库存为 $A_m$，在途采购量为 $T_m$，安全库存缺口为 $G_m$，增长系数为 $\beta = 1.05$，则采购基线为

\begin{equation}
C_m = \max(B_m, \beta H_m, G_m)
\end{equation}

最终建议采购量为

\begin{equation}
Q_m^{buy} = \max(0, C_m - A_m - T_m)
\end{equation}

该公式与当前项目中的实现逻辑一致：系统综合 BOM 需求、历史采购量、库存缺口和在途采购量后给出建议，而非只根据单一阈值触发采购。已有测试样例表明，在 BOM 需求 20、上周采购量 4、库存 2、在途 1 的条件下，系统能够计算得到建议采购量 17，并正确保留优选供应商信息。

\subsection{计划生成流程说明}

为便于概括项目实现中的周计划处理过程，本文给出算法~\ref{alg:weekly-plan-real}。

\begin{algorithm}[H]
    \KwIn{销售订单集合、生产完工记录、采购订单集合、库存记录、BOM 数据、安全库存阈值}
    \KwOut{生产周计划、采购周计划}
    统计上一周各成品完工量；
    汇总当前未完成订单对成品的需求量；
    计算当前可用成品库存；
    \ForEach{成品 $p$}{
        计算生产基线 $B_p = \max(1.10H_p, D_p)$；
        计算建议生产量 $Q_p^{prod} = \max(0, B_p - A_p)$；
        若 $Q_p^{prod} > 0$，写入生产周计划明细；
    }
    依据生产周计划与 BOM 关系换算原材料需求；
    统计上一周原材料采购量、当前可用库存与在途采购量；
    \ForEach{原材料 $m$}{
        计算采购基线 $C_m = \max(B_m, 1.05H_m, G_m)$；
        计算建议采购量 $Q_m^{buy} = \max(0, C_m - A_m - T_m)$；
        若 $Q_m^{buy} > 0$，写入采购周计划明细；
    }
    保存计划主表与明细表；
    广播生产和采购计划通知。
    \caption{周计划自动生成流程}
    \label{alg:weekly-plan-real}
\end{algorithm}

总体而言，该模块虽然不属于严格意义上的高级智能决策算法，但其工程优势非常明显：一方面它复用了系统现有数据源，另一方面它能够将经验性库存判断规则固化为可执行计算逻辑，从而增强业务管理的主动性。

\section{实时消息推送与状态同步实现}

在多角色协同系统中，仅依靠用户主动刷新页面或查询接口，难以及时获知关键业务变化。为此，本文系统在 HTTP 接口之外引入 WebSocket/STOMP 机制，将关键业务事件转换为可订阅消息，实现面向角色和业务对象的主动推送。

从实现结构看，系统通过 \texttt{WebSocketConfig} 注册 \texttt{/ws} 连接端点，并启用 \texttt{/topic} 和 \texttt{/queue} 两类消息前缀。后端各业务服务通过 \texttt{NotificationService} 或 \texttt{SimpMessagingTemplate} 将事件广播到相应主题；前端在登录成功后基于 SockJS 和 STOMP 建立连接，按当前角色订阅相关主题。例如，订单协同中会用到 \texttt{/topic/orders}、\texttt{/topic/orders/warehouse}、\texttt{/topic/orders/sales} 和按顾客邮箱归一化生成的客户主题；采购协同中会用到 \texttt{/topic/procurement}、\texttt{/topic/procurement/manager} 和按供应商账号归一化生成的供应商主题；质检模块则会用到 \texttt{/topic/quality} 以及按生产管理员邮箱生成的生产主题。

为说明其状态同步意义，可将消息推送过程抽象为事件驱动模型：设业务事件集合为 $E$，消息主题集合为 $T$，存在映射

\begin{equation}
\eta: E \rightarrow T
\end{equation}

表示每类业务事件都会投递到相应主题。当事件 $e$ 发生时，所有订阅 $\eta(e)$ 的客户端均可在近实时条件下获得通知。该机制使得系统能够将“状态变化发生”与“相关角色知晓变化”之间的时间间隔显著缩短。

项目中的典型消息场景包括：销售接受订单后通知仓库管理员；仓库发货后通知顾客；库存不足时通知生产管理员；供应商发货后通知采购管理员；质检不合格时通知对应生产管理员；质检合格时通知仓库管理员；周计划生成后通知生产管理员和采购管理员。实践表明，实时消息机制的引入使系统从被动查询式协同转变为事件驱动式协同，大幅降低了人工催办成本。

\section{质量追溯功能实现}

质量追溯功能的关键不在于单独建立一个查询页面，而在于在业务处理过程中持续保留可追溯数据。本文系统采用“订单号贯穿 + 批次承接 + 质检记录补充 + 库存流水侧证”的设计方式，使顾客和管理人员能够根据订单号追踪对应成品质检结果。

从数据实现路径看，订单对象首先提供追溯入口；生产完工时，系统将订单号写入 \texttt{Batch.sourceOrderNo}；质检处理时，\texttt{QualityRecord} 记录检验人员、检验时间、检验结果与备注；仓库入库和发货阶段则通过库存流水记录批次号和关联业务对象。因此，当顾客发起追溯请求时，后端可以先根据订单号定位批次，再进一步查询质检记录，从而给出“是否检验、是否合格、由谁检验、何时检验”等信息。

从流程保障角度看，质量追溯之所以成立，还因为系统通过质检结果反向约束业务流转。若质检不合格，订单会被退回生产环节；若质检合格，订单才能进入待生产入库。这意味着系统中的质检记录并非事后补录信息，而是真正影响发货前置条件的控制变量。正因为质量状态与订单状态、批次状态存在联动关系，顾客在查看质量追溯结果时，才能看到与实际业务一致的数据链路。

总体而言，质量追溯功能提升了系统透明度，也增强了顾客对交付质量的可验证感。对企业内部而言，该功能还能辅助识别生产和质检流程中的薄弱环节，为后续质量治理与责任分析提供客观依据。

